% chapter-2.tex
% Chapter 2 — Background and Related Work
%
% Citations use the citation keys of the project Zotero library; the
% corresponding entries were exported to references.bib. Claims that would
% require a source not currently available in the library are marked with
% % TODO comments. Pending design decisions use visible \todo{...}
% placeholders (a \providecommand fallback is included so the chapter
% compiles standalone).

\providecommand{\todo}[1]{\textbf{[TODO: #1]}}

\chapter{Background and Related Work}
\label{chap:background}

\section{Computer Vision on Embedded Platforms}
\label{sec:cv-embedded}

Computer vision is the discipline concerned with extracting information
from images and video. Over the past decade, the dominant approach to
tasks such as image classification, object detection, and segmentation has
shifted from hand-engineered features to deep neural networks, which learn
visual representations directly from data; this shift is the founding
premise of the edge deep-learning literature \cite{xu2025}. Object
detection---the task of localizing and classifying multiple objects in an
image---is a canonical perception primitive for robotic systems, and it is
commonly performed either by two-stage detectors, which first propose
candidate regions and then classify them, or by one-stage detectors, which
predict detections directly from dense feature maps and generally trade
some accuracy for higher processing speed.
% TODO: cite a dedicated survey of deep-learning-based object detection
% covering the one-stage / two-stage taxonomy; no such survey is currently
% in the Zotero library. The edge benchmarking studies cited below use
% one-stage families exclusively.
In practice, real-time embedded applications overwhelmingly adopt
one-stage detector families such as YOLO, SSD, and EfficientDet, which are
the standard subjects of edge deployment studies
\cite{alqahtani2026, szegedi2026, zequn2026}.

Deploying such models on board a robot or autonomous device imposes
constraints that do not apply to offline or data-center processing. The
perception pipeline must run continuously under real-time or near-real-time
deadlines, within a strict power and thermal envelope, and with the
computational resources of an embedded processor rather than a workstation;
inference on edge devices must balance speed, energy, memory, temperature,
and accuracy simultaneously \cite{zequn2026}. To satisfy these constraints,
practitioners typically combine three ingredients
\cite{xu2025, szegedi2026, zequn2026}. First, moderately sized or
compressed network architectures: model families are offered in graduated
sizes (e.g., nano to medium variants), and techniques such as knowledge
distillation can shrink an accurate teacher model into an edge-suitable
student, as demonstrated for medical image classification with deployment
on the Orin Nano itself \cite{aalishah2025}. Second, reduced numerical
precision: FP16 and INT8 arithmetic reduce latency, memory footprint, and
energy at a modest accuracy cost \cite{szegedi2026, zequn2026}. Third,
vendor-optimized inference runtimes that exploit the platform's GPU
accelerators; on Jetson platforms, TensorRT consistently achieves the
lowest inference latency and substantially lower energy per workload than
unoptimized frameworks \cite{szegedi2026, zequn2026}.

Independently of the inference engine, an embedded computer-vision
application has a characteristic pipeline structure: images are acquired
from a sensor or read from storage, preprocessed into the format expected
by the model, passed through the model for inference, and postprocessed to
produce application-level outputs such as bounding boxes, labels, or
annotated images. Although inference usually dominates the compute budget,
the surrounding stages are not negligible: for object-detection workloads,
preprocessing and postprocessing represent a substantial portion of total
runtime, especially at higher input resolutions or when reduced precision
accelerates inference \cite{szegedi2026}. Each stage may therefore be
computationally significant, and the way the stages are
connected---within one process or across several---is precisely the
architectural question this thesis investigates.

\section{NVIDIA Jetson Architecture}
\label{sec:jetson}

The NVIDIA Jetson family comprises embedded computing modules that
integrate Arm CPU cores, an NVIDIA GPU, and memory on a single board,
designed for robotics and edge artificial-intelligence workloads
\cite{scalcon2024}. This thesis uses the Jetson Orin Nano, a member of the
Orin generation of modules. Only hardware characteristics relevant to the
experiments are described here; the exact module configuration is
documented in Chapter~3.

The Orin Nano combines a hexa-core Arm CPU with an NVIDIA GPU of the Ampere
generation that includes Tensor Cores for the accelerated matrix
operations used in neural-network inference.
\todo{VERIFY: CPU core count, GPU architecture, and Tensor Core presence
against the official NVIDIA Jetson Orin Nano datasheet for the specific
module used.}
A characteristic of the platform that is directly relevant to this study is
its unified physical memory: the CPU and GPU access the same LPDDR5 memory
rather than separate host and device memories.
\todo{VERIFY: memory type and capacity of the specific module.}
This shared-memory design influences how expensive it is to move image
data between processing stages and between processes, which is one of the
costs introduced by distributing a pipeline across multiple processes.

Two further platform characteristics matter for reproducible benchmarking.
First, the module operates in selectable power modes that trade performance
against power consumption; published evaluations of the Orin family
explicitly distinguish these configurations---for example, the 15~W
configuration of the Orin Nano---when reporting achievable video-analytics
frame rates \cite{scalcon2024}, so the power mode must be fixed and
documented for all experiments. Second, sustained workloads are subject to
thermal effects: under continuous inference load, device temperature rises
measurably and must be monitored when characterizing performance
\cite{zequn2026}. The cooling configuration is therefore likewise fixed and
documented.
% TODO: cite the official NVIDIA Jetson Orin Nano datasheet / developer
% documentation for power modes and thermal behavior; not currently in the
% library.

The software environment for Jetson platforms is provided by JetPack, a
bundle that includes the Linux operating system, CUDA GPU computing
libraries, and accelerated inference libraries such as TensorRT
\cite{szegedi2026, zequn2026}. The exact versions used in this thesis are
recorded in Chapter~3.

\section{ROS 2 Architecture}
\label{sec:ros2}

ROS~2 is an open-source framework for developing robot software as a
distributed collection of cooperating processes \cite{macenski2022,zotero-item-9}. Its architectural elements are nodes, topics, and messages.

A \emph{node} is a process that encapsulates one coherent piece of
functionality, such as reading a camera, preprocessing images, or running
a neural-network model. Nodes exchange data by publishing and subscribing
to \emph{topics}, which are named buses carrying typed \emph{messages}. A
publisher node writes messages to a topic, and every subscriber node to
that topic receives its own copy of the message. Message types are defined
declaratively, and standard types exist for the image data used by
computer-vision applications \cite{macenski2022, zotero-item-9}.

A node executes application logic in \emph{callbacks}: functions invoked
in response to events such as the arrival of a message or the expiration
of a timer. \emph{Executors} are the mechanism that takes queued callbacks
and executes them. The ROS~2 executor model has been analyzed in detail by
the real-time-systems community \cite{casini2026}: the standard
distributions provide a single-threaded executor, which processes
callbacks sequentially within one thread, a multi-threaded executor, which
can execute callbacks concurrently, and a static single-threaded variant.
Executor choice and callback structure determine how a pipeline stage
responds to incoming data, and it is impossible to decouple the executor
from the performance of a node \cite{macenski2023, casini2026}. Executor
settings therefore constitute part of the architectural configuration
under study.

ROS~2 offers several client libraries for different programming languages.
This thesis uses \emph{rclpy}, the Python client library. Most users
interact with ROS~2 exclusively through such client libraries, which
provide the core communication APIs and rest on common lower layers: an
intermediate C library, \emph{rcl}, which provides functionality shared by
all client libraries, and below it the \emph{rmw} (ROS MiddleWare)
abstraction layer, through which the underlying communication
implementation is selected \cite{macenski2022, zotero-item-9}. This
layering matters for the interpretation of measurements: profiling work on
the C++ stack has shown that a significant share of end-to-end latency can
originate above the transport itself, in the notification and retrieval
path between the middleware and the client library
\cite{kronauer2021, kamiya2025}, so overhead observed in a Python ROS~2
pipeline may originate in the client-library and serialization layers
rather than in the transport alone.

Finally, ROS~2 supports \emph{composable nodes}: multiple nodes loaded
into a single process, which can communicate through intra-process
mechanisms that avoid serialization and copying \cite{macenski2023}.
Composition is the recommended design pattern of ROS~2 precisely because
it mitigates the costs of the multi-process default \cite{macenski2022,macenski2023}. Composable nodes are nonetheless explicitly excluded from
the scope of this thesis (Section~\ref{sec:scope}): the object of study is
the ordinary multi-process node architecture itself, and composition is
discussed as future work in Chapter~6.

\section{ROS 2 Communication}
\label{sec:ros2-communication}

At a conceptual level, ROS~2 delegates inter-process communication to an
implementation of the Data Distribution Service (DDS) standard, a
publish-subscribe middleware for distributed real-time systems
\cite{gambo2025, macenski2022}. Unlike ROS~1, whose communication depended
on a central master (\emph{rosmaster}), DDS implementations discover peers
automatically, match publishers with subscribers by topic and type, and
exchange data without any central broker \cite{macenski2022, gambo2025}.
DDS is suitable for real-time distributed embedded systems because of its
configurable transport behavior and scalability \cite{maruyama2016}, and
DDS-based middleware is now widespread in robotics, supporting applications
from multi-robot coordination to real-time sensor processing
\cite{gambo2025}.

Two aspects of this design are directly relevant to this thesis. First,
DDS exposes \emph{quality-of-service} (QoS) policies that determine
delivery guarantees and buffering behavior: \emph{reliability}
(best-effort or reliable), \emph{durability} (e.g., volatile versus
transient-local), \emph{history} settings (keep-last queues of a given
depth or keep-all), and further policies such as deadline, lifespan, and
liveliness intended for real-time system design \cite{macenski2022,zotero-item-9}. QoS configuration measurably affects latency and resource
usage: enabling reliable delivery and DDS-Security features can increase
latency substantially, and shallow history depth combined with reliable
delivery can cause retransmission failures \cite{fernandez2020,kronauer2021}. For this reason, a single fixed QoS configuration is used
throughout the experiments of this thesis. Second, the rmw abstraction
layer allows different DDS implementations to be selected at runtime
\cite{maruyama2016, macenski2022}, and the choice measurably affects
communication performance in a hardware-dependent way
\cite{maruyama2016, kronauer2021, zhang2024}. The implementation is
likewise fixed for all experiments, since comparing middlewares is outside
the scope of this thesis.

When a message travels between two nodes in separate processes, it
traverses a path that includes serialization on the publisher side,
transport over sockets or shared memory, buffering in subscriber queues,
deserialization, and finally delivery of the callback to the executor.
Each step consumes time and memory, and each has been quantified in the
literature. Profiling of the ROS~2 call stack shows that the overhead
relative to raw DDS communication depends on payload size: for small
messages, latency is dominated by the DDS middleware itself and by the
delay between DDS notifying the client library and the executor actually
retrieving the data, whereas for large payloads the DDS transport
dominates outright \cite{kronauer2021}. Layer-resolved measurements that
timestamp the boundaries between the ROS~2 and DDS layers attribute more
than half of total latency in default socket-based communication to the
DDS transport segment, and show that shared-memory transport reduces
overall latency to less than half \cite{kamiya2025}. Even intra-process
communication, which passes pointers instead of serializing, exhibits
publishing costs that grow with message size and with the number of
subscribers that require their own message copies \cite{luo2025}.
Large messages are precisely the weak point of default configurations:
UDP buffer sizes, packet fragmentation, and reliability retransmissions
combine to degrade performance for payloads on the order of megabytes
\cite{macenski2022, kronauer2021}---and the ROS~2 literature itself
identifies images and point clouds as the canonical payloads of this size
\cite{macenski2022, macenski2023}. For image messages, whose size grows
with resolution, the per-message cost of serialization and copying can be
expected to grow accordingly; whether this translates into a measurable
end-to-end penalty at the resolutions tested is an empirical question that
the experiments of this thesis are designed to answer.

\section{Computer-Vision Workload}
\label{sec:workload}

The thesis uses a single representative computer-vision workload: an
object-detection pipeline that processes a stream of images and produces
bounding-box detections. Object detection is chosen because it is a
representative and well-studied perception workload for the target
platform: published evaluations demonstrate feasible real-time object
detection on the Orin Nano with models of the YOLO, SSD, and EfficientDet
families \cite{scalcon2024, alqahtani2026, szegedi2026, zequn2026}. The
exact neural-network model will be selected according to the following
criteria:

\begin{itemize}
  \item Ability to run reliably on the Jetson Orin Nano
  \item Reproducibility
  \item Availability of documented model weights
  \item Reasonable inference time
  \item Suitability for real-time or near-real-time processing
\end{itemize}

\todo{SPECIFY: the final selected model, its source publication, weights
version, inference runtime, and any precision reduction applied. This
section must be completed once the model is chosen.}

The same model and inference configuration will be used in every benchmark
configuration, so that application architecture is the only primary
difference between experimental conditions.

\section{Performance Metrics}
\label{sec:metrics}

This section defines the performance metrics used in the thesis. The set
follows common practice in the ROS~2 performance literature, where
end-to-end latency is decomposed into communication, computation, and idle
components \cite{betz2023a} and reported together with throughput and
resource utilization \cite{mayoral-vilches2024, macenski2022}.

\subsection{End-to-End Latency}

Latency is the time between an input entering the pipeline and the
corresponding output being generated---for the pipelines of this thesis,
the elapsed time from the moment an image frame is handed to the first
processing stage until the fully postprocessed result for that same frame
is available. Latency is a distributional quantity: following the practice
of ROS~2 measurement studies \cite{kronauer2021, macenski2022}, it is
summarized by mean, median, standard deviation, minimum, maximum, and the
95th and 99th percentiles where sufficient samples are available. Tail
percentiles matter because perception results that arrive late may be as
unusable as those that never arrive \cite{betz2023a}.

\subsection{Throughput}

Throughput is the rate at which frames complete the entire pipeline,
measured in frames per second (FPS), computed as processed frames divided
by elapsed time.

\subsection{CPU Utilization}

CPU utilization is the fraction of CPU capacity in use, reported as average
and, where useful, peak utilization over the measurement interval.

\subsection{Memory Utilization}

Memory utilization is system RAM usage over the measurement interval; GPU
memory usage may also be recorded if reliable measurement is available.

\subsection{GPU Utilization}

GPU utilization records whether architectural overhead changes the
utilization of the primary compute resource for inference.

\subsection{Power Consumption (Optional)}

If reliable and reproducible measurement is possible with the available
hardware, the following quantities will be recorded: average power, peak
power, energy per processed frame, and performance per watt, mirroring the
power-related metrics used in Orin Nano characterization studies
\cite{zequn2026}. Power measurement is treated as an optional extension
rather than a requirement.
% TODO: cite Jetson power-measurement methods (e.g., on-board power rails
% / tegrastats) if power measurement is included; official documentation
% is not currently in the library.

\section{Related Work}
\label{sec:related-work}

The literature relevant to this thesis can be organized into five
categories.

\paragraph{ROS 2 communication benchmarks.}
A body of work benchmarks the communication performance of ROS~2 itself.
Maruyama et al.\ provided the earliest systematic exploration, showing
that the ROS~2 layers over DDS add overhead and that transport latency
varies with data size, configuration, and DDS vendor \cite{maruyama2016}.
Kronauer et al.\ measured end-to-end latency in multi-node pipelines of 3
to 23 nodes across three DDS middlewares on a desktop PC and a Raspberry
Pi~4, profiling the call stack and finding overheads of up to 50\% relative
to raw DDS for small messages, roughly 30\% message loss for 500~KB
payloads on the embedded device, and non-linear latency growth beyond
roughly 13--17 nodes at rates above 50~Hz \cite{kronauer2021}. Fernandez
et al.\ quantified the effect of QoS reliability settings and DDS-Security
on latency and throughput \cite{fernandez2020}. Macenski et al.\
benchmark inter-process, composed, and intra-process communication in the
ROS~2 design review: with one publisher and one subscriber, 95th-percentile
latency remains below 1~ms for inter-process messages up to 1~MB but rises
to 7.85~ms at 8~MB, with sustainable rates falling from 1000~Hz to
213~Hz \cite{macenski2022}. Kamiya et al.\ instrumented the internal
layer boundaries of shared-memory communication \cite{kamiya2025}, and
Luo et al.\ analyzed the scalability of intra-process communication,
proposing message pooling to reduce repeated allocation costs
\cite{luo2025}. These microbenchmark studies characterize the middleware
with synthetic payloads and publisher--subscriber chains, predominantly in
C++; they do not embed the middleware in an end-to-end application in
which inference dominates the workload.

\paragraph{Robotics middleware studies.}
Broader middleware research frames ROS~2 in context. Systematic reviews
map the role of DDS-based middleware across robotics applications and its
open challenges, including performance and scalability
\cite{gambo2025}, and surveys consolidate the architectural changes from
ROS~1 to ROS~2 and the ecosystem built on them \cite{al-batati2026}, as
well as the real-time support and analysis developed for ROS~2
\cite{casini2026}. Zhang et al.\ compare DDS against the alternative
Zenoh and MQTT middlewares for edge-to-edge and edge-to-cloud ROS~2
communication across Ethernet, Wi-Fi, and 4G, finding the best choice to
be network-dependent \cite{zhang2024}. Portugal et al.\ surveyed the ROS
community on the state of ROS~2 adoption, documenting both the middleware
meeting user expectations and the reasons that hold teams back from
migrating \cite{portugal2025}. These studies concern middleware selection
and adoption rather than the architecture of a single application on a
single host, which is the focus of this thesis.

\paragraph{Embedded computer-vision benchmarks.}
Another line of work benchmarks computer-vision and deep-learning
inference on embedded platforms. Alqahtani et al.\ evaluated nine
object-detection models across Raspberry Pi devices with and without TPU
accelerators, a Jetson Nano, and a Jetson Orin Nano, reporting trade-offs
among accuracy, inference time, and energy and finding the Orin Nano to
offer the most favorable overall balance \cite{alqahtani2026}. Szegedi
et al.\ compared model architectures, programming languages, inference
frameworks, and numerical precisions on an Orin Nano against a desktop
workstation \cite{szegedi2026}, and comparable characterizations exist for
power, energy, and thermal behavior under sustained inference
\cite{zequn2026}. Scalcon et al.\ assessed Orin-family devices for
AI-powered video monitoring, reporting CPU, GPU, and RAM utilization for a
surveillance workload \cite{scalcon2024}. Beyond single-model studies,
resource-management work examines GPU sharing among concurrent edge AI
applications on Jetson hardware \cite{woo2025}, and hardware-aware model
design targets efficient deployment on the Orin Nano itself
\cite{aalishah2025}. These studies establish achievable inference
performance on the target class of devices, but they consider the
application as a single process and therefore do not address the cost of
architectural partitioning.

\paragraph{Jetson platform characterization.}
Studies in this category characterize the Jetson platforms themselves:
achievable video-analytics frame rates across Orin models and power
configurations \cite{scalcon2024}, and power, energy, and temperature
behavior under sustained neural-network load \cite{zequn2026}. They
inform the experimental controls of this thesis---fixed power mode,
fixed cooling configuration, warm-up, and temperature
monitoring---rather than its research question directly.

\paragraph{ROS 2 communication overhead in applications.}
The category closest to this thesis comprises studies that investigate
the overhead ROS~2 introduces in realistic systems. Macenski et al.\
benchmarked node composition on full robotic systems, reporting 28\% CPU
and 33\% RAM savings from composing a complete navigation stack into one
process on ARM and concluding that image-scale multi-process pipelines
are virtually impossible without composition \cite{macenski2023}. Betz
et al.\ developed an evaluation workflow for the open-source autonomous
driving stack Autoware.Universe, decomposing end-to-end latency into
communication, computation, and idle components and identifying
bottlenecks across software and hardware configurations
\cite{betz2023a}, and later evaluated a containerized microservice
variant of the same stack, finding mean end-to-end latency improvements
of 5--8\% over bare-Linux deployment \cite{betz2024}. Adjacent to
robotics, the software-engineering literature on monolithic versus
microservice architectures reports that monoliths outperform their
microservice counterparts on a single machine \cite{blinowski2022}, that
the performance gap is measurable and workload-dependent
\cite{tapia2020}, and even identifies a trend of migrating back from
distributed to monolithic designs for reasons of cost, complexity, and
performance \cite{felisberto2024}. What none of these studies provides is
a comparison in which the workload, model, input, and hardware are held
equivalent while the application architecture itself is the variable
under study, with a single-process baseline, on an embedded platform, in
Python.

\section{Research Gap}
\label{sec:gap}

The reviewed literature can be summarized by what each strand leaves
unanswered. Middleware microbenchmarks characterize ROS~2 communication
in isolation, with synthetic payloads and mostly C++ clients, but not
within a complete vision pipeline in which inference dominates and frames
the communication costs \cite{maruyama2016, kronauer2021, kamiya2025,luo2025, macenski2022}. Mitigation studies---composition, intra-process
communication, shared-memory transport, message pooling---evaluate
alternatives to the multi-process architecture and thereby presuppose,
rather than measure, the cost of that architecture against a monolithic
baseline \cite{macenski2023, kamiya2025, luo2025}. Inference benchmarks
on Jetson platforms characterize the dominant computation but assume a
monolithic application and thus leave the communication dimension
unexplored \cite{szegedi2026, zequn2026, alqahtani2026, scalcon2024}.
Full-stack studies such as the Autoware evaluations measure real systems
but conflate countless algorithmic and architectural factors, without an
equivalent single-process reference \cite{betz2023a, betz2024}. The
general monolith-versus-microservices evidence, for its part, is not
robotics-specific and does not involve ROS~2, image payloads, or embedded
hardware \cite{tapia2020, blinowski2022, felisberto2024}.

What remains to be characterized---and what this thesis measures---is the
performance cost of partitioning a single, functionally equivalent Python
computer-vision pipeline across ROS~2 rclpy nodes on a single embedded
platform, as a function of the degree of partitioning, under a fixed DDS
implementation and QoS configuration, and with a reproducible methodology.
